% Nouveau EGA V (cote 25): archive cote 25, site page 81 (PDF page 82).
% typescript. Draft compared: none.
% AI-transcribed by Claude Opus 5.5 (claude-opus-5-5, adaptive thinking, effort high), Message Batches API, batch ega25b (msgbatch_01A7n5YuXubUPNK1zsepi8RG), 2026-11-13; not yet checked by eye.
\scanpage{25/81}
\mspage{\hand{76}}

NB \quad Le raisonnement ne marche plus si on remplace x par une
sous-variété linéaire \unsure{$X$|$C$} de dim.$>0$ , et qu'on assujettit H à passer
par C: en effet, rien ne s'oppose (prenant par exemple dim C = r-2, le
plus grand possible pour que H puisse encore effectivement varier) que
T contienne la droite $C^{o}$ :\xout{(et pour faire un exemple où où ce fait se}
\xout{produit, en car 0 disons, on \unsure{part} d'une sous-variété T de} prendre par
exemple une quadrique \inserted{non singulière} dans $P$ , en caractéristique quelconque. Mais il
est possible que de tels phénomènes ne puissent non plus se produire \xout{xx}
pour des $\phi_n f$ , $n\geq 2$ ; on pourrait poser la question en remarque, au
nº 9 .

\uline{Remarques} 13.4. \quad a) Nous avons déjà signalé que l'hypothèse que k \xout{xx}
soit parfait est essentielle pour la validité de 13.2. Par contre, il
est plausible que l'hypothèse ``k infini'' est surabondante. \xout{En tous cas}
Nous n'avons pas essayé \xout{xxx} de faire un raisonnement ad hoc pour le cas
où k est fini, et signalons seulement que dans ce cas, l'application de
13.2. à la clôture algébrique de k et les raisonnements habituels
montrent qu'on peut trouver une extension finie k' de k telle que
\xout{xx} \inserted{pour} les points de $X_{k'}$ au dessus de x il existe des voisinages ouverts
polyélementaires relativement à k'.

b) Lorsque dans 13.2. on abandonne l'hypothèse que X est géom
irréductible, la conclusion bien entendu ne subsiste pas telle quelle,
(puisqu'un schéma algébrique poly-élementaire est géom irréd !). Elle
subsiste cependant sous la forme plus faible qu'on obtiendrait en
omettant dans la définition 13.1. le mot ``connexes'' : c'est ce que
montre la démonstration que nous avons donnée.
\inserted{c)}

\over{c}{d}) En procédant comme expliqué dans 7.1. on peut \inserted{donner} des variantes
de 13.2. dans le cas où on remplace le corps de base k par un préschéma
\margin{81}
